\subsection{The source's columnwise harmonic cancellation}
\label{subsec:sr-source-harmonic-chart-cancellation}

Fix the same packet, coefficient ring and framed point
as in Subsection~\ref{subsec:sr-source-frame-truncation}.
A letter \(\lambda\in\mathrm{Lt1}(D_w)\) has an
underlying embedding \(\tau_\lambda\); its multiplicity
index is retained.
For a multi-index \(\alpha\) write
\(\alpha!=\prod_\lambda\alpha_\lambda!\).
The exact exponential-coordinate value is
\[
 \mathrm{expVal}(g,\alpha)=
      \alpha!\,[z^\alpha]g .
 \tag{HC1}
\]
For \(n\ge0\), rename the variables of \(g_b\) into
column \(b\) and put
\[
 {\cal G}(z)=\prod_{b=1}^n g_b(z_b),\qquad
 \langle P,F\rangle=P(\partial)F(0).
\]
The coordinate bijection is precisely
\((\lambda,b)\leftrightarrow(\tau_\lambda,i_\lambda,b)\)
of the actual source.  It gives
\[
 \langle P,{\cal G}\rangle
 =\sum_{\mu\in\mathrm{supp}P}P_\mu
                 \prod_b\mathrm{expVal}(g_b,\mu_b).
 \tag{HC2}
\]
In this and subsequent formulas coefficients of \(P\)
are mapped from \(c.F\) to \(S\).
There is no factorial for the number of columns:
(HC2) uses exactly the product of the coordinate
factorials.

\begin{srlemma}[The exact finite product-rule cancellation]
\label{lem:sr-source-quadratic-product-rule-cancellation}
Let \(S\) be any commutative \(\mathbb Q\)-algebra,
\(Q\) a homogeneous quadratic polynomial, and let
\(c_{ij}=\partial_i\partial_jQ\).
Put
\[
                   {\cal D}_cP=\sum_{i,j}c_{ij}
                                          \partial_i\partial_jP.
\]
For every polynomial \(P\) and formal series \(A\),
\[
                  \langle P,AQ\rangle
                         =\tfrac12\langle{\cal D}_cP,A\rangle .
 \tag{HC3}
\]
If \({\cal D}_cP=0\), then for every \(t\in S\)
and every formal series \(F\),
\[
                         \langle P,e^{tQ}F\rangle
                                      =\langle P,F\rangle. \tag{HC4}
\]
\end{srlemma}
\begin{proof}
Expand the ordered derivatives defining a monomial
\(P(\partial)\) on \(AQ\) by the product rule.
Terms with zero or one derivative applied to \(Q\)
vanish at the origin because \(Q(0)=dQ(0)=0\).
Terms with more than two vanish because \(Q\) is
quadratic.
Choose the two derivative positions applied to \(Q\).
The remaining derivatives act on \(A\), and those two
positions contribute \(c_{ij}\).
The operator \({\cal D}_c\) counts the two orders of
each chosen pair; symmetry of \(c\) makes their
contributions equal.
Division by \(2\) gives (HC3), including repeated
derivative indices.
Extend by coefficient linearity to every \(P\).
This is the source's \(\mathrm{half\_gen}\) ordered-word
argument, with its factor \(1/2\).

To obtain (HC4), use the source's factorization
\[
 e^{tQ}=1+tQY,\qquad
                  Y=\sum_{r\ge0}\frac{(tQ)^r}{(r+1)!}.
 \tag{HC5}
\]
All substitutions are legitimate formal substitutions
since \(Q(0)=0\); only finitely many terms affect any
specified coefficient.
Apply (HC3) to \(A=tFY\).
The extra term in \(\langle P,F(1+tQY)\rangle\)
is zero because \({\cal D}_cP=0\).
Thus no analytic convergence or division by a
nonzero coefficient of \(P\) is used.
\end{proof}

For \(\beta=(\beta_{\lambda\lambda'})\), put
\[
 Q_\beta^{(n)}
       =\sum_{b=1}^n\sum_{\lambda,\lambda'}
                     \beta_{\lambda\lambda'}
                                    z_{\lambda,b}z_{\lambda',b}.
 \tag{HC6}
\]
The matrix \(\beta\) need not be symmetric.
Its Hessian is
\[
 c_{(\lambda,b),(\lambda',b')}
  =\mathbf1_{b=b'}(\beta_{\lambda\lambda'}
                                      +\beta_{\lambda'\lambda}).
 \tag{HC7}
\]
This includes both terms and their exact common-column
condition.

\begin{srproposition}[Cancellation for the literal harmonic space]
\label{prop:sr-source-actual-column-harmonic-cancellation}
Let \(P\in\mathrm{harmA}(c,a,n)\), and suppose
\[
 \beta_{\lambda\lambda'}\ne0
               \ \Longrightarrow\
              \tau_{\lambda'}=\overline{\tau_\lambda}.
 \tag{HC8}
\]
Then for every \(t\in S\),
\[
 \sum_\mu P_\mu\prod_b
        \mathrm{expVal}(g_be^{tQ_\beta^{(1)}},\mu_b)
       =\sum_\mu P_\mu\prod_b\mathrm{expVal}(g_b,\mu_b).
 \tag{HC9}
\]
The same equality holds if each new column series
is specified only modulo terms of degree \(D+1\),
provided every monomial of \(P\) has degree at most \(D\).
\end{srproposition}
\begin{proof}
We first identify the exact harmonic input.
The definition of \(\mathrm{harmA}\), through its
identity-matrix Laplacian, gives
\[
 \sum_b\partial_{(\tau,i,b)}
                  \partial_{(\bar\tau,j,b)}P=0
                                  \qquad(\tau\in\Psi_H).   \tag{HC10}
\]
Every conjugate pair has a member in \(\Psi_H\).
Indeed the host element \(\mathrm{rt}\) has square
\(-d\) with \(d>0\); the chosen complex square root
\(\mathrm{rt}_0\) is nonzero and purely imaginary,
so \(\overline{\mathrm{rt}_0}=-\mathrm{rt}_0\).
Any embedding sends \(\mathrm{rt}\) to
\(\mathrm{rt}_0\) or \(-\mathrm{rt}_0\).
In the second case its conjugate sends it to
\(\mathrm{rt}_0\).
This is the literal \(\mathrm{PsiCover}\) argument
for the actual host, rather than an extra selection
of embeddings.
If the first embedding in a conjugate pair is not
in \(\Psi_H\), use (HC10) in the reverse order and
commute the two polynomial partial derivatives.

Rename the source letters by the displayed coordinate
bijection and map coefficients to \(S\).
Renaming by an injection and coefficient extension
commute with polynomial partial derivatives.
Thus for every pair satisfying (HC8),
\[
              \sum_b\partial_{\lambda,b}
                              \partial_{\lambda',b}P=0.   \tag{HC11}
\]
Use (HC7).
For each ordered pair, the term involving
\(\beta_{\lambda\lambda'}\) is zero by (HC11) if
that coefficient is nonzero, and is zero directly
otherwise.
The term involving \(\beta_{\lambda'\lambda}\)
is zero by the reversed (HC11).
Hence \({\cal D}_cP=0\).

The product of the column exponential factors is
\(e^{tQ_\beta^{(n)}}\), by (HC6) and the formal
exponential addition law.
Apply (HC4) to \({\cal G}=\prod_b g_b(z_b)\), and
use (HC2) on both sides.
This proves (HC9) with the source's factorials.

For the finite-order assertion, substitution into a
column and multiplication of finitely many series
preserve equality modulo \((z)^{D+1}\).
The pairing with \(P\) reads only coefficients in
degrees at most \(D\), so it cannot detect any of
those residual terms.
This proves the exact finite-jet version as well.
For \(n=0\), the products are empty products and the
same proof gives the constant case.
\end{proof}

We now apply this calculation to an actual chart's
formal group and canonical gauge, before making any
comparison with a second chart.
Let \(\ell_G\) be the actual normalized logarithm in
(FC2), let \(\varepsilon_G\) be its compositional
inverse, and set
\[
 U(z)=\sum_\lambda u_\lambda z_\lambda,\qquad
                          \phi(z)=\varepsilon_G(U(z)).
 \tag{HC12}
\]
The inverse exists by the source's recursive formal
inverse construction: the linear term of \(\ell_G\)
is \(I\), and after the inverse has been determined
through degree \(r-1\), the degree-\(r\) error is
cancelled by its negative, since the unknown
degree-\(r\) term enters with coefficient \(I\).
The construction gives the two compositional
inverse identities coefficientwise.
It works over the same arbitrary \(S\).

Thus \(\phi(0)=0\), \(d\phi(0)e_\lambda=u_\lambda\),
and
\[
              G(\phi(z),\phi(w))=\phi(z+w),\qquad
                   A_a\phi(z)=\phi(D_a z).                \tag{HC13}
\]
The first follows by applying \(\ell_G\) and then
its inverse.
For the second, (FC2) and (FC15) give
\(\ell_G(A_a\phi(z))=A_a^{(1)}U(z)=U(D_a z)\);
apply the inverse again.
These are the exact \(\mathrm{IsExpOf}\) and
action-equivariance clauses used in the source.

Let \(D_{G,u}\) denote the source invariant derivative,
whose coefficient vector is \(J(X)u\).
Differentiate the first identity in (HC13) with respect
to \(w_\lambda\) at \(w=0\) to get
\[
                  \partial_\lambda\phi(z)
                                      =J(\phi(z))u_\lambda.
\]
The formal chain rule therefore gives, for every \(h\),
\[
          \partial_\lambda(h\circ\phi)
                              =(D_{G,u_\lambda}h)\circ\phi.
 \tag{HC14}
\]
Iterating this identity in the source's ordered
derivative word gives the same assertion for each
monomial and then by linearity for a polynomial.
Substitution by the zero-constant \(\phi\) preserves
the constant coefficient.
Consequently
\[
 (D_{G,u}^{\alpha}h)(0)
                        =\mathrm{expVal}(h\circ\phi,\alpha).
 \tag{HC15}
\]
The factorial in (HC15) follows by successively using
\([z^\nu]\partial_i h=(\nu_i+1)[z^{\nu+e_i}]h\).
This is the actual \(\mathrm{tie\_alg}\) calculation.

For a chart \(C\) with its supplied canonical gauge
\(h\), write \(h_L(X)=C.B.hL(h,w_0)\), exactly as
defined by the source graph \((X,A_{w_0}X)\).
The definition of \(\mathrm{nabVal}\) and (HC15)
then give
\[
 \begin{split}
 g_b&=(f_b e^{kh_L})\circ\phi,\\
 \mathrm{chartVal}(C,k,(f_b),P)
      &=\sum_\mu P_\mu
                       \prod_b\mathrm{expVal}(g_b,\mu_b).
 \end{split}                                                   \tag{HC16}
\]
This is a formal identity for the actual supplied
canonical gauge; at this point it asserts no new
existence of exponential line generators.

\begin{srproposition}[Invariance under the source frame correction]
\label{prop:sr-source-actual-harmonic-frame-correction}
Suppose \(B^t=B\) is balanced as in (FC9), and
\[
       j_Df'_b=j_D(f_b e^{kQ_B})\quad\text{for every }b.
\]
For \(P\in\mathrm{harmA}(c,a,n)\) with all support
degrees at most \(D\),
\[
             \mathrm{chartVal}(C,k,(f'_b),P)
                         =\mathrm{chartVal}(C,k,(f_b),P).
 \tag{HC17}
\]
This supplies the literal \(\mathrm{HarmInvSt}\).
\end{srproposition}
\begin{proof}
From (HC12)--(HC13) and (FC13),
\[
 Q_B(\phi(z))
   =\sum_{\lambda,\lambda'}
       \left(u_\lambda^tBu_{\lambda'}\right)
                \chi_{\lambda'}(w_0)z_\lambda z_{\lambda'}.
 \tag{HC18}
\]
By (FC16) the displayed coefficient matrix satisfies
(HC8).
Substitution by \(\phi(0)=0\) preserves the given
degree-\(D\) jet equality.
Multiplication by the common \(e^{kh_L}\) and formal
exponential substitution give
\[
 j_D\bigl((f'_be^{kh_L})\circ\phi\bigr)
        =j_D\bigl(g_b e^{kQ_B\circ\phi}\bigr).
\]
Apply the finite-order form of (HC9) with \(t=k\)
and the matrix in (HC18), and use (HC16).
This proves (HC17), with no assertion that the
original pair generators already agree under
truncation.
Together with (FC14), it removes the correction
from each actual harmonic chart value.
\end{proof}

The same cancellation with \(t=-k\), and with the
quadratic \(q\) of (GR5), removes the ratio (GR6)
between two actual exponential gauges on one actual
exponential chart.
For precise coefficients, write
\[
 q(z)=\sum_{\lambda\le\lambda'}
            [z_\lambda z_{\lambda'}]q\,
                                   z_\lambda z_{\lambda'}.
\]
Use the upper-triangular coefficient matrix in (HC6);
there is no duplication or extra factor \(2\) here.
Its Hessian is still exactly (HC7).
It has (HC8) by (GR5).
This is the source's literal
\(\mathrm{KillAddValSt}\), and its finite-order
version used in \(\mathrm{HarmInvSt}\).

\paragraph{Frozen source correspondence.}
With prefix
\texttt{src/ReciprocityTheory/ReductiveGroups/ArithmeticGeometry/}
\texttt{TaylorDiscriminants/WickVanishing/},
(HC3) is \(\mathrm{half\_gen}\) in
\texttt{FoldrIterateEqTG.lean}
(\texttt{f78ffd8df71057d7bbfc2d572e72d45e2f125abc});
its ordered-word splitting supplies the repeated-index
count and factor \(1/2\).
(HC5) and the resulting \(\mathrm{kill\_gen}\)
application are \texttt{KillAdd.lean}
(\texttt{9d2613f7dec448fbda02d08b983f7ae7fd93265b}).
The exact Hessian (HC7) is
\texttt{KillAddQ.lean}
(\texttt{21c18e91c9c2c5753e554b27c9a618253f4cc029}).
\texttt{KillAddFin.lean}
(\texttt{dd02e3e4ef65cda13260d786dd84b767e2c6d1d6})
is the assembly \(\mathrm{kernel\_additive}\);
the product-rule argument above supplies its underlying
proof rather than using that assembly alone.

All following files have prefix
\texttt{src/Formalism/SharpRestriction/Divisibility/GeneratorPowers/}.
The exact Laplacian (HC10) and coordinate reindexing are
\texttt{CIndEncPf.lean}
(\texttt{52edd9a81392677693c49162dc23274853586166}).
The actual host \(\Psi\)-coverage is
\texttt{PsiCover.lean}
(\texttt{7c2eff14be2863c5d6ba777d3f7bf12ae034b667}).
The coefficient-factorial and product-of-columns
identities (HC1)--(HC2), and the explicit Hessian
annihilation, are \texttt{ConstantCoeffPolyAdd.lean}
(\texttt{a7bd8d585f47122e5ec3f71e38ffc06d9a0d90cc}).
The upper-triangular quadratic coefficient matrix and
ordinary exponential-gauge cancellation are
\texttt{KillAddVal.lean}
(\texttt{8400f516d63d26349f192b3d83a8ac8e3c839614}).

The actual formal exponential (HC12) is
\texttt{ExpOfPf.lean}
(\texttt{f664907a8e698159c83b34e49a504ac516332ee6});
the action is \(\mathrm{expActOf}\) using (FC15).
The chain-rule, ordered-iterate and factorial calculation
(HC14)--(HC16) is \texttt{ExpTieAlg.lean}
(\texttt{7edcde1ac562ca6d5944a80930274029ad278b79}),
\(\mathrm{tie\_alg}\), \(\mathrm{nabVal\_eq\_const}\).
The logarithm substitution (HC18) is
\texttt{HarmInvB.lean}
(\texttt{32c2974b0108fe6014d111c901644035585898c8});
the finite-jet product step is \texttt{HarmInvC.lean}
(\texttt{c6552267ccc6f92d2cafbae69323022bd7a87832});
the actual \(\mathrm{HarmInvSt}\) assembly is
\texttt{HarmInvPf.lean}
(\texttt{31abf027367853bef3b186f174937d035375b5e8}).

This section proves the two harmonic cancellation
suppliers and the formal exponential derivative
identity.  The actual exponential-gauge construction
and the equality of actual exponential chart point
systems remain the two geometric comparisons required
before arbitrary-chart independence is concluded.
